\uuid{Jl2n}
\exo7id{799}
\titre{exo7 799}
\auteur{cousquer}
\organisation{exo7}
\datecreate{2003-10-01}
\isIndication{false}
\isCorrection{false}
\chapitre{Calcul d'intégrales}
\sousChapitre{Théorie}
\cpge{mpsi-mp2i.fonctions-valeurs-reelles-complexes.b, pcsi.techniques-calcul-fonctions.b, ptsi.techniques-calcul-fonctions.b, mpsi-mp2i.primitives-equations-differentielles.a, pcsi.techniques-calcul-primitives.a, ptsi.techniques-calcul-primitives.a, mpsi-mp2i.integration.e, pcsi.integration.d, ptsi.integration.d, bcpst1.limites-continuite.c, bcpst1.integration-segment.c, mp-mpi.fonctions-vectorielles.c}
\module{Analyse}
\niveau{L1}
\difficulte{}

\contenu{
\texte{
Soit $\varphi$ une fonction bornée sur $[a,b]$~; comparer les
assertions suivantes\footnote%
    {L'une des implications à étudier est très difficile~; on pourra 
    admettre après avoir traité toutes les autres que celle qui reste 
    est fausse.}~:
}
\begin{enumerate}
    \item \question{$\varphi$ a une primitive sur $[a,b]$.}
    \item \question{$\varphi$ est intégrable sur $[a,b]$.}
    \item \question{$\varphi$ est continue sur $[a,b]$.}
    \item \question{$\varphi$ est dérivable sur $[a,b]$.}
\end{enumerate}
}
